% =============================================================================
% Textos para el libro — revisión P2 (comparación estadística del OE4)
% Cifras: resultados_fase4/reservados/resumen_estadistico.json
%         (sensibilidad: resultados_fase4/completo/resumen_estadistico.json)
% Generado con: python evaluar_fase4_estadistica.py [--fuente completo]
% =============================================================================


% -----------------------------------------------------------------------------
% 1. METODOLOGÍA — Fase 4 (reemplaza la descripción de la cola única)
% -----------------------------------------------------------------------------

\paragraph{Comparación estadística.}
El agente se compara con las seis políticas de referencia (Solo CPU, Solo GPU, MET,
MCT, Min-Min y Max-Min) sobre 100 colas de 200 tareas. Las colas se generan a partir
de los tamaños reservados en la prueba de interpolación (Sección~\ref{sec:interpolacion}),
es decir, con tareas cuyo tamaño de problema el agente no observó durante el
entrenamiento: la cola de semilla $s \in \{0, \dots, 99\}$ se muestrea sin reemplazo
entre las tareas reservadas de la partición $k = s \bmod 5$, de modo que cada partición
aporta 20 colas, y se evalúa con los cinco agentes entrenados sin esa partición. El EDP
de sistema del agente en una cola es la media de esos cinco agentes. Todas las políticas
se evalúan sobre las mismas colas, con el tiempo de sistema definido como el
\textit{makespan} de las colas paralelas de CPU y GPU y el EDP de sistema como el producto
de la energía total por ese tiempo.

Para cada cola $i$ y cada política $P$ se calcula la diferencia relativa de EDP de sistema
respecto del agente,
\begin{equation}
  d_{i}(P) = 100 \left( \frac{\mathrm{EDP}_{i}(P)}{\mathrm{EDP}_{i}(\text{agente})} - 1 \right) \%,
\end{equation}
de modo que $d_{i}(P) > 0$ indica que el agente obtiene menor EDP que la política $P$ en
esa cola. Para cada una de las seis comparaciones se aplica la prueba de rangos con signo
de Wilcoxon pareada por cola ($H_0$: la mediana de $d(P)$ es cero), con aproximación normal
y corrección por empates, y los seis valores $p$ se ajustan con la corrección de Holm
($\alpha = 0{,}05$). El tamaño del efecto se reporta como $r = |Z|/\sqrt{n}$, con $n$ el
número de colas con diferencia no nula. Se reportan la mediana y el rango intercuartílico
(IQR) de $d(P)$.

Antes de calcular las pruebas se fijó un umbral de relevancia práctica de 1\,\% del EDP de
sistema: con 100 colas pareadas, diferencias del orden de 0,1\,\% pueden resultar
estadísticamente significativas sin tener importancia práctica. Una diferencia se considera
relevante solo si es significativa tras la corrección de Holm y su mediana supera el umbral
en valor absoluto. El oráculo por tarea (el dispositivo de menor EDP para cada tarea, conocido
a posteriori) se reporta solo como referencia descriptiva, porque no es una política
implementable. Como análisis de sensibilidad se repite el procedimiento con colas muestreadas
del conjunto completo y el modelo final del agente.


% -----------------------------------------------------------------------------
% 2. RESULTADOS — Fase 4: nueva subsección
%    (reemplaza la tabla tab:fase4_interpolacion de una cola por partición)
% -----------------------------------------------------------------------------

\subsection{Comparación estadística sobre 100 colas}
\label{sec:fase4_estadistica}

La Figura~\ref{fig:fase4_robustez} muestra la distribución, cola a cola, de la diferencia de
EDP de sistema de cada política respecto del agente, y la Tabla~\ref{tab:fase4_wilcoxon}
resume las pruebas. Las seis diferencias son significativas tras la corrección de Holm
($p < 0{,}001$), pero solo cuatro superan el umbral de relevancia práctica.

\begin{figure}[H]
  \centering
  \includegraphics[width=\textwidth]{figuras/fase4_robustez.png}
  \caption{Diferencia de EDP de sistema de cada política respecto del agente en 100 colas
  de 200 tareas de tamaños no observados. Cada punto es una cola; la caja marca el IQR y
  la línea gruesa la mediana. La franja gris corresponde al umbral de relevancia de
  $\pm 1$\,\%. Solo CPU se muestra en un panel aparte por su escala.}
  \label{fig:fase4_robustez}
\end{figure}

\input{tablas/tabla_fase4_wilcoxon.tex}

El agente supera con diferencias relevantes a Solo CPU (mediana $+1397$\,\%, en las 100
colas) y a Max-Min (mediana $+3{,}74$\,\%, IQR $[+1{,}44;\ +5{,}43]$\,\%, mejor en 89 de 100
colas). Frente a Solo GPU ($+0{,}23$\,\%) y MET ($+0{,}06$\,\%) las diferencias son
significativas pero quedan por debajo del umbral: en la práctica, el agente es equivalente a
estas dos políticas. MCT y Min-Min obtienen menor EDP de sistema que el agente, con
diferencias relevantes: medianas de $-4{,}34$\,\% (IQR $[-7{,}48;\ -2{,}88]$\,\%) y
$-3{,}34$\,\% (IQR $[-4{,}76;\ -2{,}05]$\,\%), y el agente solo las supera en 3 y 4 de las 100
colas, respectivamente.

El agente queda a una mediana de 0,02\,\% del oráculo por tarea, es decir, resuelve casi sin
error el problema para el que fue entrenado: elegir, para cada tarea, el dispositivo de menor
EDP. La desventaja frente a MCT y Min-Min no proviene de errores de clasificación, sino de la
formulación del problema. El estado del agente describe solo la tarea actual, no la ocupación
de las colas, y su recompensa es inmediata ($\gamma = 0$). Por tanto, no puede equilibrar el
\textit{makespan} entre CPU y GPU. MCT y Min-Min, en cambio, envían a la CPU parte de las
tareas en que la GPU es más rápida. En la mediana por cola, frente al agente, reducen el
\textit{makespan} de 16,2~s a 13,0~s y 13,7~s, a costa de 15\,\% y 7\,\% más energía. Como el
EDP de sistema multiplica la energía total por el \textit{makespan}, ese equilibrio compensa.
Max-Min reduce el tiempo en una proporción similar (13,3~s), pero con 27\,\% más energía, y
termina con mayor EDP que el agente.

Los resultados no dependen de la semilla del agente: al usar por separado cada uno de los
cinco agentes de cada partición, las medianas cambian como máximo 0,01 puntos porcentuales.
Con colas tomadas del conjunto completo y el modelo final del agente se obtienen las mismas
conclusiones (Solo GPU $+0{,}21$\,\%, MET $+0{,}04$\,\%, MCT $-4{,}95$\,\%, Min-Min
$-2{,}94$\,\%, Max-Min $+3{,}91$\,\%; todas con $p_{\mathrm{Holm}} < 0{,}01$).


% -----------------------------------------------------------------------------
% 3. CONCLUSIONES — OE4 (reemplaza el texto actual)
% -----------------------------------------------------------------------------

Sobre 100 colas de 200 tareas de tamaños no observados, comparadas con la prueba de Wilcoxon
pareada y corrección de Holm, el agente redujo el EDP de sistema de forma significativa y
relevante frente a Solo CPU (mediana 1397\,\%) y frente a Max-Min (mediana 3,7\,\%), y fue
equivalente en la práctica a Solo GPU y MET (diferencias de 0,23\,\% y 0,06\,\%, por debajo
del umbral de 1\,\%). MCT y Min-Min obtuvieron un EDP de sistema menor que el del agente en
4,3\,\% y 3,3\,\% (medianas). El agente alcanza casi el óptimo por tarea (0,02\,\% del
oráculo), pero una política que decide tarea por tarea, sin observar la ocupación de las
colas, no captura el beneficio de equilibrar el \textit{makespan} que aprovechan MCT y
Min-Min. El OE4 se cumple como comparación: la evaluación cuantifica con significancia
estadística y relevancia práctica en qué casos el agente mejora a las políticas de referencia
y en cuáles no.


% -----------------------------------------------------------------------------
% 4. CONCLUSIÓN GENERAL — frase a incorporar
% -----------------------------------------------------------------------------

En la evaluación de sistema sobre 100 colas de tamaños no observados, el agente redujo el EDP
frente a Solo CPU y Max-Min (medianas de 1397\,\% y 3,7\,\%), igualó en la práctica a Solo GPU
y MET, y quedó por detrás de MCT y Min-Min (medianas de $-4{,}3$\,\% y $-3{,}3$\,\%), que
equilibran el \textit{makespan} entre dispositivos; incorporar la ocupación de las colas al
estado del agente es la extensión natural para cerrar esa brecha.


% -----------------------------------------------------------------------------
% 5. RESUMEN — frase a incorporar (sustituye las cifras de la cola única)
% -----------------------------------------------------------------------------

Sobre 100 colas de 200 tareas de tamaños no observados (prueba de Wilcoxon pareada con
corrección de Holm y umbral de relevancia de 1\,\%), el agente redujo el EDP de sistema frente
a Solo CPU y Max-Min (medianas de 1397\,\% y 3,7\,\%), fue equivalente a Solo GPU y MET, y
obtuvo un EDP 4,3\,\% y 3,3\,\% mayor que MCT y Min-Min.


% -----------------------------------------------------------------------------
% 6. LIMITACIONES
%    - RETIRAR: «Evaluación sobre una única cola».
%    - AGREGAR (opcional, sustenta la discusión de MCT/Min-Min):
% -----------------------------------------------------------------------------

\textbf{Decisión por tarea, sin estado de las colas.} El estado del agente describe la tarea
actual y su recompensa es inmediata; no incluye la ocupación de las colas de CPU y GPU. Por
ello el agente optimiza el EDP de cada tarea, pero no el \textit{makespan} del lote, que es lo
que aprovechan MCT y Min-Min en el EDP de sistema. Incorporar esa información al estado queda
como trabajo futuro.
